% !TEX root = ../rl_book.tex
\chapter{Introduction}

Reinforcement Learning (RL) is both a problem and solution framework
%studies how an \emph{agent} can learn effective behavior through trial and error while interacting with an \emph{environment}. The agent repeatedly observes, acts, and receives scalar feedback (rewards), and the objective is to learn a behavior that maximizes cumulative reward (the \emph{return}). Deep RL couples these ideas with function approximation via deep networks to represent policies and value functions. \cite{spinningup_rl_intro}

\begin{figure}[h]
\centering
\begin{tikzpicture}[
    >=Stealth,
    looparrow/.style={->, line width=0.9pt, draw=gray!70},
    agentnode/.style={draw, rounded corners=6pt, fill=teal!15, minimum width=3.8cm, minimum height=1.2cm},
    envnode/.style={draw, rounded corners=6pt, fill=blue!10, minimum width=3.8cm, minimum height=1.2cm},
    font=\sffamily
]
  % Nodes
  \node[agentnode] (agent) {Agent};
  \node[envnode, below=10mm of agent] (env) {Environment};

  % Curved arrows with labels
  \draw[looparrow] (agent.east) to [out=0, in=0, looseness=1.25]
    node[midway, right=6mm, align=center] {Action\\$a_t$} (env.east);

  \draw[looparrow] (env.west) to [out=180, in=180, looseness=1.25]
    node[midway, left=7mm, align=center] {State, Reward\\$s_t,\, r_t$} (agent.west);
\end{tikzpicture}

\medskip
\emph{Agent–environment interaction loop.}
\end{figure}


\section{What Can RL Do?}
RL has achieved notable results in game playing (e.g., Go, Dota, Atari), continuous control for robotics (both simulated and real), and instruction-following tasks. These successes illustrate the breadth of domains where learning through interaction and delayed feedback is beneficial. \cite{spinningup_rl_intro}

\section{Key Concepts and Terminology}

\subsection{Agent--Environment Loop}
At discrete time $t=0,1,2,\dots$, the agent receives an observation $o_t$ (often a function of an underlying state $s_t$), selects an action $a_t$, and the environment transitions to $s_{t+1}$ while emitting a reward $r_t$. The goal is to choose actions so as to maximize a cumulative notion of return. \cite{spinningup_rl_intro}

\subsection{States and Observations}
A \emph{state} $s_t$ is a complete description of the world; an \emph{observation} $o_t$ may reveal only part of that state. In deep RL, states/observations are typically represented as real-valued vectors, matrices, or tensors (e.g., pixels for vision, joint angles/velocities for robotics). \cite{spinningup_rl_intro}

\subsection{Action Spaces}
The \emph{action space} may be \emph{discrete} (finite set of actions) or \emph{continuous} (real-valued vectors). This distinction strongly influences algorithm design: many methods are tailored either for discrete or for continuous control. \cite{spinningup_rl_intro}

\subsection{Policies}
A \emph{policy} maps observations or states to actions. Deterministic policies are often written
\[
a_t = \mu_\theta(s_t),
\]
while stochastic policies are written
\[
a_t \sim \pi_\theta(\cdot \mid s_t),
\]
with parameters $\theta$ (e.g., neural network weights) learned by optimization. In practice, categorical policies are common for discrete actions and (diagonal) Gaussian policies for continuous actions. \cite{spinningup_rl_intro}

\subsection{Trajectories}
A \emph{trajectory} (episode) is $\tau=(s_0,a_0,r_0,s_1,a_1,r_1,\ldots)$, generated by rolling out the policy in the environment with its transition dynamics. \cite{spinningup_rl_intro}

\subsection{Reward and Return}
Two standard notions of return are:
\begin{align}
\text{Finite-horizon, undiscounted:}\quad & R(\tau)=\sum_{t=0}^{T} r_t, \\
\text{Infinite-horizon, discounted:}\quad & R(\tau)=\sum_{t=0}^{\infty} \gamma^{t} r_t,\qquad \gamma\in(0,1).
\end{align}
Discounting is both intuitive (earlier rewards count more) and mathematically convenient (guarantees convergence of infinite sums under mild conditions). \cite{spinningup_rl_intro}

\subsection{The RL Objective}
For a policy $\pi_\theta$, the objective is to maximize expected return:
\[
\max_{\theta}\; J(\theta) \;=\; \mathbb{E}_{\tau\sim\pi_\theta}\!\left[\,R(\tau)\,\right].
\]
The expectation is taken over trajectories induced by $\pi_\theta$ and the environment dynamics. \cite{spinningup_rl_intro}

\subsection{Value Functions}
Value functions quantify expected return under a policy $\pi$:
\begin{align}
V^{\pi}(s) &= \mathbb{E}\!\left[ \sum_{t=0}^{\infty}\gamma^t r_t \;\middle|\; s_0=s,\; a_t\sim \pi(\cdot\mid s_t) \right], \\
Q^{\pi}(s,a) &= \mathbb{E}\!\left[ \sum_{t=0}^{\infty}\gamma^t r_t \;\middle|\; s_0=s,\; a_0=a,\; a_{t\ge1}\sim \pi(\cdot\mid s_t) \right].
\end{align}
The corresponding optimal value functions are
\[
V^{*}(s)=\max_{\pi} V^{\pi}(s),\qquad Q^{*}(s,a)=\max_{\pi} Q^{\pi}(s,a).
\]
\cite{spinningup_rl_intro}

There are two key connections between the value function and the action-value function:

\[
V^\pi(s) = \mathbb{E}_{a \sim \pi}[\,Q^\pi(s,a)\,], 
\qquad
V^*(s) = \max_{a} Q^*(s,a).
\]

These follow directly from the definitions:
\begin{itemize}
    \item $V^\pi(s)$ is the expected return starting from $s$ while following $\pi$. Since the first action $a$ is chosen from $\pi(\cdot|s)$, this expectation equals the average of $Q^\pi(s,a)$ under $\pi$.
    \item $V^*(s)$ represents the optimal return achievable from $s$. This is simply the value of the best possible action at $s$, i.e.\ $\max_a Q^*(s,a)$.
\end{itemize}

In the infinite-horizon discounted case, the value function is
\[
V^\pi(s) = \mathbb{E}_\pi \!\left[\sum_{t=0}^\infty \gamma^t r_t \,\Big|\, s_0=s \right],
\]
which depends only on the state $s$ since the problem is stationary and the discount ensures convergence.  

For a finite-horizon undiscounted return
\[
R_t = \sum_{k=t}^{T-1} r_k, \qquad 
V_t^\pi(s) = \mathbb{E}_\pi \!\left[\sum_{k=t}^{T-1} r_k \,\Big|\, s_t=s \right],
\]
the value must depend on time $t$: the same state can have different values depending on how many steps remain. In particular, $V_T^\pi(s)=0$ since no rewards are left once the horizon is reached. Hence, unlike the infinite-horizon discounted case, the finite-horizon setting breaks stationarity and requires $V_t^\pi(s)$.

The optimal action-value function $Q^*(s,a)$ gives the expected return for starting
in state $s$, taking action $a$, and then following the optimal policy thereafter.
The optimal policy selects the action that maximizes this return:
\[
a^*(s) = \arg\max_a Q^*(s,a).
\]
If multiple actions achieve the maximum, any of them is optimal; thus there may be
several optimal policies, but at least one deterministic optimal policy always exists.



\subsection{Bellman Equations}
Value functions satisfy self-consistency (Bellman) relations. For on-policy values:
\begin{align}
V^{\pi}(s) &= \mathbb{E}_{a\sim \pi(\cdot\mid s),\, s'\sim P(\cdot\mid s,a)}\!\left[\, r(s,a) + \gamma V^{\pi}(s') \right], \\
Q^{\pi}(s,a) &= \mathbb{E}_{s'\sim P(\cdot\mid s,a),\, a'\sim \pi(\cdot\mid s')}\!\left[\, r(s,a) + \gamma Q^{\pi}(s',a') \right].
\end{align}
For optimal values:
\begin{align}
V^{*}(s) &= \max_{a}\; \mathbb{E}_{s'\sim P(\cdot\mid s,a)}\!\left[\, r(s,a) + \gamma V^{*}(s') \right], \\
Q^{*}(s,a) &= \mathbb{E}_{s'\sim P(\cdot\mid s,a)}\!\left[\, r(s,a) + \gamma \max_{a'} Q^{*}(s',a') \right].
\end{align}
The right-hand sides are sometimes called \emph{Bellman backups}. \cite{spinningup_rl_intro}

\subsection{Relationship between $V^\pi$ and $Q^\pi$}

The state-value and action-value functions under a policy $\pi$ are related by the following expectation identities:

\paragraph{Expectation over actions (under the policy).}
Since the state-value is the expected return given state $s$, averaging over actions according to the policy,
\begin{equation}
    V^\pi(s) 
    = \mathbb{E}_{a \sim \pi(\cdot|s)}\big[ Q^\pi(s,a) \big]
    = \sum_{a \in \mathcal{A}} \pi(a|s)\, Q^\pi(s,a).
\end{equation}
Here the expectation is taken with respect to the action distribution $\pi(\cdot|s)$.

\paragraph{Expectation over next states (under environment dynamics).}
The action-value is defined as the expected return given state $s$ and action $a$. Using the Bellman equation,
\begin{equation}
    Q^\pi(s,a) 
    = \mathbb{E}_{s' \sim P(\cdot|s,a)}\big[ r(s,a) + \gamma V^\pi(s') \big],
\end{equation}
where the expectation is taken with respect to the next-state distribution 
$P(s'|s,a)$ induced by the environment dynamics.


\subsection{Advantage Functions}
The \emph{advantage} of taking $a$ in $s$ under $\pi$ is
\[
A^{\pi}(s,a) \;=\; Q^{\pi}(s,a) - V^{\pi}(s),
\]
which measures how much better the action is than the policy’s average choice at that state. Advantages are central to many policy-gradient estimators. \cite{spinningup_rl_intro}

\section{(Optional) Formalism: Markov Decision Processes}
A Markov Decision Process (MDP) is a tuple $\langle S,A,R,P,\rho_0\rangle$ with state space $S$, action space $A$, reward function $R:S\times A\times S\to\mathbb{R}$, transition kernel $P:S\times A\to\mathcal{P}(S)$, and initial-state distribution $\rho_0$. The Markov property asserts that $s_{t+1}$ depends only on $(s_t,a_t)$, not on the full past. \cite{spinningup_rl_intro}

\bigskip
\noindent\textbf{Notation.} Throughout, expectations $\mathbb{E}[\cdot]$ are taken over the relevant randomness from the policy and environment; when needed, we use $s'\sim P(\cdot\mid s,a)$ and $a\sim \pi(\cdot\mid s)$ to emphasize sampling.


\begin{figure}[h]
\centering
\begin{tikzpicture}[
  >=Stealth,
  box/.style={draw, rounded corners=6pt, minimum width=4.2cm, minimum height=1.2cm, font=\sffamily, align=center},
  samp/.style={box, fill=orange!15},
  val/.style={box,  fill=blue!12},
  pol/.style={box,  fill=green!15},
  lab/.style={font=\sffamily\footnotesize, align=center}
]

% Nodes arranged in a triangle
\node[samp] (sampling) {Sampling};
\node[val, below right=25mm and -10mm of sampling] (value) {Value Evaluation};
\node[pol, below left=25mm and -10mm of sampling] (policy) {Policy Improvement};

% Curved clockwise arrows
\draw[->, thick] (sampling.east) to [out=0,in=90] node[lab, right=2mm] {trajectories $\mathcal{D}$} (value.north);
\draw[->, thick] (value.south west) to [out=240,in=300] node[lab, below=2mm] {critic / gradients} (policy.south east);
\draw[->, thick] (policy.north) to [out=120,in=180] node[lab, left=2mm] {policy $\pi_\theta$} (sampling.west);

% Notes below each box
\node[lab, below=2mm of sampling] {Collect rollouts with $\pi_\theta$};
\node[lab, below=2mm of value] {Estimate $V^\pi$, $Q^\pi$, $\hat A_t$};
\node[lab, below=2mm of policy] {Update $\theta$ (e.g., PG, PPO, TRPO)};

\end{tikzpicture}

\medskip
\caption{Typical RL algorithm pipeline.}
\end{figure}


\section{Summary}
In value-based reinforcement learning, the error of a $Q$-function update is defined as an 
expectation over the true transition dynamics:
\[
Q(s,a) \;\leftarrow\; r + \gamma \, \mathbb{E}_{s' \sim P(\cdot|s,a)} \big[ \max_{a'} Q(s',a') \big],
\]
where the expectation is taken over the actual environment transition distribution 
$P(s'|s,a)$. 

In contrast, policy-gradient methods shift the focus to the distribution of actions 
induced by the policy:
\[
\nabla_\theta J(\pi_\theta) \;=\; \mathbb{E}_{s \sim d^{\pi}, \, a \sim \pi_\theta(\cdot|s)}
\!\left[ \nabla_\theta \log \pi_\theta(a|s)\, Q^{\pi}(s,a) \right],
\]
where the expectation depends on the state--action distribution under the current 
policy $\pi_\theta$. 

Thus, value-based methods rely directly on transition samples, while 
policy-based methods rely on the action distribution defined by the policy itself.
